%! The Secant, Cosecant, and Cotangent Functions — Unit 7, Lesson 7.4 — Guided Notes (Answer Key)
\documentclass[10pt]{article}
\usepackage{precalculus-article}
\usepackage{precalculus-key}   % pulls in -boxes; do NOT also load -boxes
\newcommand{\TallMath}[1]{$\displaystyle #1\rule[-1.4em]{0pt}{3.2em}$}

\begin{document}

\pageheader{Unit 7, Lesson 7.4}{Guided Notes --- Answer Key}
\namedateperiod

\vspace{0.10in}

%----------------------------------------------------------------------
% Objectives
%----------------------------------------------------------------------
\begin{objectivebox}
\textbf{I will be able to\ldots}
\begin{enumerate}[leftmargin=1.5em, itemsep=4pt]
  \item \textbf{LO 3.11.A} --- Define secant, cosecant, and cotangent as
    reciprocals of cosine, sine, and tangent, respectively, and evaluate
    each at unit-circle angles. \writeline
  \item \textbf{LO 3.11.A} --- Identify the domain, range, vertical asymptotes,
    and period of $\sec\theta$, $\csc\theta$, and $\cot\theta$. \writeline
\end{enumerate}
\end{objectivebox}

\vspace{0.08in}

%----------------------------------------------------------------------
% Vocabulary
%----------------------------------------------------------------------
\begin{vocabbox}
\begin{description}[leftmargin=0pt, itemsep=6pt]
  \item[\termblanklong{secant}]
    $\sec\theta = \dfrac{1}{\ans{$\cos\theta$}}$;
    the reciprocal of cosine.
    Undefined when $\cos\theta = $ \ans{0}.

  \item[\termblanklong{cosecant}]
    $\csc\theta = \dfrac{1}{\ans{$\sin\theta$}}$;
    the reciprocal of sine.
    Undefined when $\sin\theta = $ \ans{0}.

  \item[\termblanklong{cotangent}]
    $\cot\theta = \dfrac{\ans{$\cos\theta$}}{\ans{$\sin\theta$}}$;
    the reciprocal of tangent.
    Undefined when $\sin\theta = $ \ans{0}.

  \item[\termblanklong{vertical asymptote}]
    A vertical line $\theta = a$ where the function value grows without
    \ans{bound}; caused by a \ans{zero} in the denominator.

  \item[\termblanklong{reciprocal function}]
    A function formed by $\dfrac{1}{f(\theta)}$.
    Zeros of $f$ become \ans{asymptotes} of the reciprocal,
    and $f(\theta) = \pm 1$ are \ans{invariant} points (outputs unchanged).
\end{description}
\end{vocabbox}

\vspace{0.08in}

%----------------------------------------------------------------------
% Hook
%----------------------------------------------------------------------
\begin{hookbox}
\textbf{Entry Question:} The graph shown is $y = \sec\theta$ over $[0, 2\pi]$.
There are gaps where the function is undefined.

\vspace{4pt}
\textbf{Q1.} At approximately which $\theta$ values are the gaps located?
$\theta = $ \ans{$\pi/2$ and $3\pi/2$}

\vspace{4pt}
\textbf{Q2.} The output values $-0.8$ and $0.3$ never appear. What does that tell
you about the range of $y = \sec\theta$?
\ans{The range excludes the open interval $(-1, 1)$; outputs are always $\leq -1$ or $\geq 1$.}

\vspace{4pt}
\textbf{Q3.} What function do you think causes those gaps?
\ans{The cosine function: $\sec\theta$ is undefined wherever $\cos\theta = 0$.}
\end{hookbox}

\vspace{0.08in}

%----------------------------------------------------------------------
% LO 3.11.A (Part 1) — Secant and Cosecant
%----------------------------------------------------------------------
\begin{notesbox}{LO 3.11.A --- Part 1: Secant and Cosecant Functions}

\textbf{Secant:} $\sec\theta = \dfrac{1}{\cos\theta}$, where $\cos\theta \neq 0$.

\vspace{6pt}
\renewcommand{\arraystretch}{1.7}
\begin{tabular}{|c|c|c|}
\hline
$\theta$ & $\cos\theta$ & $\sec\theta = \dfrac{1}{\cos\theta}$ \\
\hline
$0$ & $1$ & \ans{1} \\
\hline
$\dfrac{\pi}{3}$ & $\dfrac{1}{2}$ & \ans{2} \\
\hline
$\dfrac{\pi}{2}$ & $0$ & \ans{undefined} \\
\hline
$\pi$ & $-1$ & \ans{$-1$} \\
\hline
$\dfrac{3\pi}{2}$ & $0$ & \ans{undefined} \\
\hline
$2\pi$ & $1$ & \ans{1} \\
\hline
\end{tabular}

\vspace{6pt}
Vertical asymptotes of $\sec\theta$: $\theta = $ \ans{$\pi/2$ and $3\pi/2$} (in $[0,2\pi]$)

\vspace{4pt}
General formula for asymptotes: $\theta = \ans{$\dfrac{\pi}{2} + n\pi$}$ for integer $n$

\vspace{4pt}
Range of $\sec\theta$: \ans{$(-\infty, -1] \cup [1, \infty)$}

\vspace{10pt}
\textbf{Cosecant:} $\csc\theta = \dfrac{1}{\sin\theta}$, where $\sin\theta \neq 0$.

\vspace{6pt}
\begin{tabular}{|c|c|c|}
\hline
$\theta$ & $\sin\theta$ & $\csc\theta = \dfrac{1}{\sin\theta}$ \\
\hline
$0$ & $0$ & \ans{undefined} \\
\hline
$\dfrac{\pi}{6}$ & $\dfrac{1}{2}$ & \ans{2} \\
\hline
$\dfrac{\pi}{2}$ & $1$ & \ans{1} \\
\hline
$\pi$ & $0$ & \ans{undefined} \\
\hline
$\dfrac{3\pi}{2}$ & $-1$ & \ans{$-1$} \\
\hline
$2\pi$ & $0$ & \ans{undefined} \\
\hline
\end{tabular}

\vspace{6pt}
Vertical asymptotes of $\csc\theta$: $\theta = $ \ans{$0$, $\pi$, and $2\pi$} (in $[0,2\pi]$)

\vspace{4pt}
General formula for asymptotes: $\theta = \ans{$n\pi$}$ for integer $n$

\vspace{4pt}
Range of $\csc\theta$: \ans{$(-\infty, -1] \cup [1, \infty)$}

\end{notesbox}

\vspace{0.08in}

%----------------------------------------------------------------------
% LO 3.11.A (Part 2) — Cotangent
%----------------------------------------------------------------------
\begin{notesbox}{LO 3.11.A --- Part 2: Cotangent Function}

\textbf{Cotangent:}
$\cot\theta = \dfrac{\cos\theta}{\sin\theta} = \dfrac{1}{\tan\theta}$
\quad (where $\sin\theta \neq 0$)

\vspace{6pt}
\renewcommand{\arraystretch}{1.7}
\begin{tabular}{|c|c|c|c|}
\hline
$\theta$ & $\sin\theta$ & $\cos\theta$ & $\cot\theta = \dfrac{\cos\theta}{\sin\theta}$ \\
\hline
$0^+$ (approaching 0) & $\approx 0^+$ & $1$ & \ans{$+\infty$} \\
\hline
$\dfrac{\pi}{4}$ & $\dfrac{\sqrt{2}}{2}$ & $\dfrac{\sqrt{2}}{2}$ & \ans{1} \\
\hline
$\dfrac{\pi}{2}$ & $1$ & $0$ & \ans{0} \\
\hline
$\dfrac{3\pi}{4}$ & $\dfrac{\sqrt{2}}{2}$ & $-\dfrac{\sqrt{2}}{2}$ & \ans{$-1$} \\
\hline
$\pi^-$ (approaching $\pi$) & $\approx 0^+$ & $-1$ & \ans{$-\infty$} \\
\hline
\end{tabular}

\vspace{6pt}
As $\theta$ increases from $0$ to $\pi$, cotangent is \ans{decreasing}
(circle: \textit{increasing} / \textit{decreasing}).

\vspace{4pt}
Vertical asymptotes of $\cot\theta$: $\theta = \ans{$n\pi$}$ for integer $n$

\vspace{4pt}
Period of $\cot\theta$: \ans{$\pi$}
\hspace{1.5cm}
Range of $\cot\theta$: \ans{$(-\infty, \infty)$}

\end{notesbox}

\vspace{0.08in}

%----------------------------------------------------------------------
% Summary Table
%----------------------------------------------------------------------
\begin{notesbox}{LO 3.11.A --- Summary: The Six Trigonometric Functions}

\renewcommand{\arraystretch}{1.8}
\begin{tabularx}{\linewidth}{@{} l X X X @{}}
\hline
\textbf{Function} & \textbf{Asymptotes} ($\theta$ excluded from domain) & \textbf{Range} & \textbf{Period} \\
\hline
$\sec\theta = 1/\cos\theta$ & \ans{$\theta = \pi/2 + n\pi$} & \ans{$(-\infty,-1]\cup[1,\infty)$} & \ans{$2\pi$} \\
$\csc\theta = 1/\sin\theta$ & \ans{$\theta = n\pi$} & \ans{$(-\infty,-1]\cup[1,\infty)$} & \ans{$2\pi$} \\
$\cot\theta = \cos\theta/\sin\theta$ & \ans{$\theta = n\pi$} & \ans{$(-\infty,\infty)$} & \ans{$\pi$} \\
\hline
\end{tabularx}

\end{notesbox}

\vspace{0.08in}

%----------------------------------------------------------------------
% Practice
%----------------------------------------------------------------------
\begin{practicebox}

\textbf{Guided Practice 1.}\quad Evaluate each reciprocal trig function.
Use the unit circle values given.

\vspace{4pt}
\renewcommand{\arraystretch}{1.7}
\begin{tabular}{|l|c|c|}
\hline
\textbf{Given} & \textbf{Compute} & \textbf{Value} \\
\hline
$\cos\!\left(\dfrac{\pi}{4}\right) = \dfrac{\sqrt{2}}{2}$ &
$\sec\!\left(\dfrac{\pi}{4}\right) = \dfrac{1}{\cos(\pi/4)} = $ & \ans{$\sqrt{2}$} \\
\hline
$\sin\!\left(\dfrac{\pi}{3}\right) = \dfrac{\sqrt{3}}{2}$ &
$\csc\!\left(\dfrac{\pi}{3}\right) = $ & \ans{$\dfrac{2\sqrt{3}}{3}$} \\
\hline
$\sin\!\left(\dfrac{\pi}{4}\right) = \dfrac{\sqrt{2}}{2}$, $\cos\!\left(\dfrac{\pi}{4}\right) = \dfrac{\sqrt{2}}{2}$ &
$\cot\!\left(\dfrac{\pi}{4}\right) = \dfrac{\cos(\pi/4)}{\sin(\pi/4)} = $ & \ans{1} \\
\hline
\end{tabular}

\vspace{0.10in}
\textbf{Guided Practice 2.}\quad For each function, identify all vertical asymptotes
in $[0, 2\pi]$ and state whether the given value is in the domain.

\vspace{4pt}
(a)\quad $f(\theta) = \sec\theta$. \quad Is $\theta = \pi/3$ in the domain?
\ans{Yes} \quad Asymptotes in $[0,2\pi]$: \ans{$\theta = \pi/2$ and $\theta = 3\pi/2$}

\vspace{6pt}
(b)\quad $g(\theta) = \csc\theta$. \quad Is $\theta = \pi$ in the domain?
\ans{No} \quad Asymptotes in $[0,2\pi]$: \ans{$\theta = 0$, $\theta = \pi$, $\theta = 2\pi$}

\vspace{6pt}
(c)\quad $h(\theta) = \cot\theta$. \quad Is $\theta = \pi/2$ in the domain?
\ans{Yes} \quad Asymptotes in $[0,2\pi]$: \ans{$\theta = 0$, $\theta = \pi$, $\theta = 2\pi$}

\end{practicebox}

\end{document}
